% !TeX program = xelatex
% Overleaf: Settings > Compiler > XeLaTeX.
% Missing screenshots render placeholders; add the real images before submission.
\documentclass[UTF8,fontset=none,11pt,a4paper]{ctexart}

% Use Traditional Chinese fonts on Overleaf, with a Windows preview fallback.
\IfFontExistsTF{Noto Serif CJK TC}{
  \setCJKmainfont{Noto Serif CJK TC}
  \setCJKsansfont{Noto Sans CJK TC}
  \setCJKmonofont{Noto Sans Mono CJK TC}
}{
  \setCJKmainfont[AutoFakeBold=2]{MingLiU}
  \setCJKsansfont{Microsoft JhengHei}
  \setCJKmonofont{MingLiU}
}

\usepackage[margin=23mm]{geometry}
\usepackage{amsmath}
\usepackage{booktabs}
\usepackage{tabularx}
\usepackage{graphicx}
\usepackage{float}
\usepackage[labelformat=empty]{caption}
\usepackage{hyperref}
\hypersetup{
  colorlinks=true,
  linkcolor=blue,
  citecolor=blue,
  urlcolor=blue,
  pdftitle={Lab 1: SGEMM Optimization},
  pdfauthor={314551147}
}
\setlength{\parindent}{0pt}
\setlength{\parskip}{0.45em}
\setlength{\emergencystretch}{2em}
\renewcommand{\arraystretch}{1.15}
\ctexset{
  section={format=\Large\bfseries},
  subsection={format=\large\bfseries}
}

% TODO: Replace only the name below. The student ID is already filled in.
\newcommand{\studentname}{楊承恩}
\newcommand{\studentid}{314551147}

% Optional graphic options can trim display margins without changing the image.
\newcommand{\reportscreenshot}[4][]{%
  \begin{figure}[H]
    \centering
    \IfFileExists{#2}{%
      \includegraphics[#1,width=\linewidth,height=0.38\textheight,keepaspectratio]{#2}%
    }{%
      \fbox{\begin{minipage}[c][18mm][c]{0.92\linewidth}
        \centering
        待補實測截圖（目前尚未加入）\\[0.4em]
        {\small\texttt{\detokenize{#2}}}
      \end{minipage}}%
    }
    \caption{#3}
    \label{#4}
  \end{figure}
}

\title{\textbf{Lab 1: SGEMM Optimization}\\[0.3em]
       \large Edge AI 實驗報告}
\author{姓名：\studentname\qquad 學號：\studentid}
\date{}

\begin{document}
\maketitle

\section*{Part 0: Know Your GPU}

\subsection*{Q1. GPU 規格與理論峰值}

\begin{center}
\begin{tabularx}{\linewidth}{@{}lX@{}}
\toprule
項目 & 數值 \\
\midrule
GPU 名稱 & NVIDIA Tesla T4 \\
Compute capability & 7.5 \\
SM 數量 & 40 \\
CUDA cores 數量 & 2560 \\
最高 Boost clock & $1590\ \mathrm{MHz}=1.59\ \mathrm{GHz}$ \\
記憶體資料傳輸速率 & 約 $10\ \mathrm{Gbps/pin}$ \\
記憶體匯流排寬度 & $256\ \mathrm{bits}$ \\
FP32 理論峰值 & $8140.8\ \mathrm{GFLOP/s}$（約 $8.14\ \mathrm{TFLOP/s}$）\\
DRAM 理論頻寬 & $320\ \mathrm{GB/s}$ \\
\bottomrule
\end{tabularx}
\end{center}

T4 的官方規格列出 2560 個 CUDA cores 與約 $8.1\ \mathrm{TFLOP/s}$
的 FP32 效能。T4 採用 Turing 架構，compute capability 為 7.5；每個 SM 含 64 個
FP32 cores，因此 $40\times64=2560$，與 CUDA cores 總數一致。

\textbf{FP32 理論峰值。}
一次融合乘加（fused multiply-add, FMA）包含一次乘法與一次加法，
計為兩次浮點運算。依作業指定的公式，以最高 Boost clock 計算：
\begin{align*}
P_{\mathrm{peak}}
  &=2\times N_{\mathrm{cores}}\times f_{\mathrm{boost}}\\
  &=2\times2560\times1.59\\
  &=8140.8\ \mathrm{GFLOP/s}.
\end{align*}
其中時脈使用 GHz，因此結果為 GFLOP/s。此結果與官方約
$8.1\ \mathrm{TFLOP/s}$ 的規格相符；它是理論上限，並非 kernel 的實測效能。

\textbf{DRAM 理論頻寬。}
以每條資料線的傳輸速率乘上匯流排寬度，再以每 byte 等於 8 bits 轉換：
\begin{align*}
B_{\mathrm{peak}}
  &=\frac{\text{data rate}\times\text{bus width}}{8}\\
  &=\frac{10\times256}{8}\\
  &=320\ \mathrm{GB/s}.
\end{align*}
此結果與官方 Product Brief 的最高 $320\ \mathrm{GB/s}$ 相符。
這裡的頻寬是 GPU 與其顯示記憶體之間的頻寬，不是主機與 GPU 之間的
PCIe 傳輸頻寬。

\subsection*{Q2. Ridge point 與 SGEMM 運算強度}

\textbf{Ridge point。}
Roofline 模型以計算峰值及記憶體頻寬界定效能上限。
兩條上限相交時的運算強度為：
\begin{align*}
I_{\mathrm{ridge}}
  &=\frac{P_{\mathrm{peak}}}{B_{\mathrm{peak}}}\\
  &=\frac{8140.8\ \mathrm{GFLOP/s}}{320\ \mathrm{GB/s}}\\
  &=25.44\ \mathrm{FLOP/byte}.
\end{align*}

\textbf{SGEMM 的運算量。}
令 $M=N=K=4096$。按照作業的計數方式，每個 C 元素需執行 $K$ 次乘加，
每次乘加計為 2 FLOP，因此主要矩陣乘法的運算量為：
\begin{align*}
F&=2MNK\\
 &=2\times4096^3\\
 &=137{,}438{,}953{,}472\ \mathrm{FLOP}.
\end{align*}

\textbf{DRAM 資料傳輸量。}
依題目假設，A 與 B 各從 DRAM 讀取一次，C 寫回一次。
每個 FP32 元素占 4 bytes；此模型不另外計入讀取舊 C 的流量。
\begin{align*}
D&=4(MK+KN+MN)\\
 &=4\times3\times4096^2\\
 &=201{,}326{,}592\ \mathrm{bytes}
  =192\ \mathrm{MiB}.
\end{align*}

\textbf{Arithmetic intensity。}
運算強度是運算量與 DRAM 資料傳輸量之比：
\begin{align*}
I_{\mathrm{SGEMM}}
  &=\frac{F}{D}
   =\frac{2MNK}{4(MK+KN+MN)}\\
  &=\frac{2\times4096^3}{4\times3\times4096^2}
   =\frac{4096}{6}\\
  &\approx682.67\ \mathrm{FLOP/byte}.
\end{align*}

\textbf{瓶頸判斷。}
由於
\[
I_{\mathrm{SGEMM}}\approx682.67
>I_{\mathrm{ridge}}=25.44\ \mathrm{FLOP/byte},
\]
此大小的 SGEMM 在題目指定的理想資料讀寫假設下屬於
\textbf{compute-bound}。每搬動一份資料所需執行的運算量很高，
因此理論上先受到 FP32 計算能力上限限制，而非 DRAM 頻寬上限。

此結論不代表 naive kernel 實測時一定為 compute-bound。
若資料未有效重複使用、產生額外 DRAM 流量，或存在其他執行限制，
實際瓶頸仍需透過後續 profiler 指標判斷。

\clearpage
\section*{Part 1: The Naive Kernel and the Profiler}

\subsection*{Q1. Benchmark 結果}

在 $M=N=K\in\{1024,2048,4096\}$、$\alpha=1$、$\beta=0$ 下，
naive 與 cuBLAS 的結果皆為 \texttt{PASS}。
Naive 的效能約為 $61\ \mathrm{GFLOP/s}$，明顯低於 cuBLAS，
且僅達 GPU FP32 理論峰值的 $0.7\%$--$0.8\%$。

\begin{table}[H]
\centering
\small
\begin{tabular}{@{}rlrrrrc@{}}
\toprule
$M=N=K$ & Kernel & ms & GFLOP/s & \%cuBLAS & \%peakF & Check \\
\midrule
1024 & naive  & 35.4488   & 60.6   & 1.6   & 0.7  & PASS \\
1024 & cuBLAS & 0.7373    & 2912.7 & 100.0 & 35.8 & PASS \\
2048 & naive  & 278.6551  & 61.7   & 1.5   & 0.8  & PASS \\
2048 & cuBLAS & 6.0286    & 2849.7 & 100.0 & 35.0 & PASS \\
4096 & naive  & 2232.1255 & 61.6   & 1.5   & 0.8  & PASS \\
4096 & cuBLAS & 34.4986   & 3983.9 & 100.0 & 48.9 & PASS \\
\bottomrule
\end{tabular}
\caption{Naive 與 cuBLAS 的 benchmark 結果。}
\label{tab:part1-benchmark}
\end{table}

Naive 在三種尺寸下約為 $60.6$--$61.7\ \mathrm{GFLOP/s}$；尺寸加倍時，
運算量與執行時間都約增加八倍。以相同尺寸的量測時間相除，
cuBLAS 分別約快 $48.08\times$、$46.22\times$ 與 $64.70\times$。
\texttt{\%cuBLAS} 以各筆 benchmark 對應的 cuBLAS 參考時間為基準，
因此不一定與兩筆獨立量測的時間比值完全相同。

\reportscreenshot{graph/1024.png}
  {$1024^3$ 的 naive 與 cuBLAS benchmark 原始輸出。}
  {fig:part1-benchmark-1024}
\reportscreenshot{graph/2048.png}
  {$2048^3$ 的 naive 與 cuBLAS benchmark 原始輸出。}
  {fig:part1-benchmark-2048}
\reportscreenshot{graph/4096.png}
  {$4096^3$ 的 naive 與 cuBLAS benchmark 原始輸出。}
  {fig:part1-benchmark-4096}

\clearpage
\subsection*{Q2. Profiler 指標、瓶頸與 FP32 峰值比例}

以下分析 $2048^3$ 的 naive kernel，從記憶體階層吞吐、
sector 存取效率與 warp stall 判斷主要瓶頸。

\textbf{GPU Speed Of Light 與記憶體階層。}
Memory Throughput 為 $49.77\%$，
比 Compute (SM) Throughput 的 $6.03\%$ 更接近 $100\%$。
這支持先檢查記憶體存取效率，而不是認為 FP32 計算管線已飽和。
不過，這裡的 Memory Throughput 是整體記憶體處理的指標，
\emph{不能當成 DRAM 頻寬使用率}。

\begin{table}[H]
\centering
\begin{tabular}{@{}lr@{}}
\toprule
GPU Speed Of Light 指標 & 實測值 \\
\midrule
Duration & $278.23\ \mathrm{ms}$ \\
Compute (SM) Throughput & $6.03\%$ \\
Memory Throughput & $49.77\%$ \\
L1/TEX Cache Throughput & $99.53\%$ \\
L2 Cache Throughput & $0.70\%$ \\
DRAM Throughput & $2.49\%$ \\
\bottomrule
\end{tabular}
\caption{Naive 在 $2048^3$ 的 GPU Speed Of Light 指標。}
\label{tab:part1-sol}
\end{table}

\reportscreenshot{graph/sol.png}
  {GPU Speed Of Light Throughput 的實測指標。}
  {fig:part1-sol}

\begin{table}[H]
\centering
\begin{tabular}{@{}lr@{}}
\toprule
Memory Workload Analysis 指標 & 實測值 \\
\midrule
Memory Throughput（DRAM 讀寫合計）& $8.00\ \mathrm{GB/s}$ \\
L1/TEX Hit Rate & $99.04\%$ \\
L2 Hit Rate & $82.07\%$ \\
\bottomrule
\end{tabular}
\caption{實際 DRAM 吞吐與 cache 命中率。}
\label{tab:part1-memory}
\end{table}

L1/TEX 的吞吐指標很高，但 L2 與 DRAM 的吞吐很低；實際 DRAM 讀寫合計
僅 $8.00\ \mathrm{GB/s}$，約為 $320\ \mathrm{GB/s}$ 理論頻寬的 $2.5\%$。
不同硬體指標的 active/elapsed 正規化基準不完全相同，不能直接相加，
或將它們解讀成資料流量的分配比例。

L1 的 sector 命中率為 $99.04\%$，表示大多數被查詢的 sector 已在 L1，
不必每次都從 DRAM 重新取得；L2 命中率則針對到達 L2 的查詢，
與 L1 命中率的分母不同。高命中率不代表每個請求的處理成本低，
因此還需要檢查 sector 數與 warp stall。

\reportscreenshot{graph/memory.png}
  {Memory Workload Analysis 的 DRAM 吞吐與 L1/L2 命中率。}
  {fig:part1-memory}

\textbf{Global-load sectors per request。}
L1/TEX 的 global-load sectors/request 為 $16.50$，
global-store sectors/request 為 $32.00$。一個 sector 為 32 bytes，
因此對 32 個 thread 各存取一個連續且對齊的 FP32 元素而言，只需 4 個 sector。

\begin{table}[H]
\centering
\small
\begin{tabular}{@{}lrrr@{}}
\toprule
存取類型 & Requests & Sectors & Sectors/Request \\
\midrule
Global Load & $537{,}001{,}984$ & $8{,}862{,}535{,}049$ & $16.50$ \\
Global Store & $131{,}072$ & $4{,}194{,}304$ & $32.00$ \\
\bottomrule
\end{tabular}
\caption{L1/TEX 的 global load/store 存取統計。比例以 sectors 除以 requests。}
\label{tab:part1-sectors}
\end{table}

Naive 的 \texttt{threadIdx.x} 對應 row；block 為 $(32,32)$，
所以同一個 warp 的 32 個 thread 負責不同 row、相同 col。
在固定的 $k$ 下，\texttt{A[row*K+k]} 在相鄰 thread 間相隔
$2048\times4=8192$ bytes，涉及 32 個不同 sector；
\texttt{B[k*N+col]} 則是所有 thread 讀同一個元素，只涉及 1 個 sector。
主要迴圈兩種讀取的平均為 $(32+1)/2=16.5$，與實測 $16.50$ 相符。
整份統計還包含迴圈外的讀取，因此實測比例不必恰好等於 $16.5$。
C 的寫回同樣跨不同列，與實測每個 store request 涉及 32 個 sector 一致。

這是\emph{同時存取的地址分散}，不是矩陣本身未連續儲存。
此外，L1 的 sector 存取數不能直接當成 DRAM 搬運量：
例如同一 thread 在相鄰 $k$ 迴圈讀取的 A 元素是連續的，
只要 sector 尚未被替換，就能在後續迴圈重用 cache 中的資料。

\reportscreenshot{graph/global.png}
  {L1/TEX global load/store 的 requests、sectors 與存取效率分析。}
  {fig:part1-global-access}

\textbf{主要 warp stall。}
LG Throttle 為 $365.25$，遠高於其他等待原因。
依 Warp State Statistics 的定義，下列值是每個 issued instruction
對應的平均 warp cycles，不是指令總數。

\begin{table}[H]
\centering
\begin{tabular}{@{}lr@{}}
\toprule
指標 & 平均週期／issued instruction \\
\midrule
Warp Cycles Per Issued Instruction & $373.61$ \\
Stall LG Throttle & $365.25$ \\
Stall Long Scoreboard & $3.41$ \\
Stall Not Selected & $2.47$ \\
\bottomrule
\end{tabular}
\caption{Warp State Statistics 的主要等待指標。}
\label{tab:part1-stalls}
\end{table}

LG Throttle 表示 local/global 記憶體指令的 L1 佇列已滿，
warp 必須等候佇列有空位才能發出下一個相關指令。
它占平均指令間隔的比例為
\[
\frac{365.25}{373.61}\times100\%\approx97.76\%.
\]
這不是整張 GPU 完全閒置的時間比例，也不是保證可達成的加速幅度。
相較之下，等待 L1/TEX 操作結果的 Long Scoreboard 僅為 $3.41$。
這些指標顯示主要壓力在記憶體指令的處理，而非單純等待 DRAM 回傳資料。
上述指標定義可參考 NVIDIA 的
\href{https://docs.nvidia.com/nsight-compute/ProfilingGuide/index.html}{Nsight Compute Profiling Guide}。

\reportscreenshot{graph/warp.png}
  {Warp State Statistics 與 LG Throttle 的實測分析。}
  {fig:part1-warp}

\textbf{達到多少 FP32 理論峰值？}
以 $2048^3$ 的一般 benchmark 結果 $61.7\ \mathrm{GFLOP/s}$，
以及 Part 0 的規格峰值 $8140.8\ \mathrm{GFLOP/s}$ 計算：
\[
\%P_{\mathrm{FP32}}
 =\frac{61.7}{8140.8}\times100\%
 \approx0.76\%.
\]
因此 harness 顯示一位小數時為 $0.8\%$。
其餘兩個尺寸也只有約 $0.74\%$--$0.76\%$。
這個比例是實際 FP32 GFLOP/s 除以規格峰值，
與 ncu 的 Compute (SM) Throughput $6.03\%$ 不同；
後者是硬體管線的吞吐指標，不能直接當成 FP32 峰值比例。

\textbf{結論。}
綜合低 FP32 效能、較高的整體記憶體處理吞吐、分散的 sector 存取，
以及主導的 LG Throttle，naive 的主要限制指向\textbf{記憶體存取的處理效率}，
而不是 FP32 math 吞吐已達上限。
DRAM 頻寬並未飽和，因此不能將此結果簡化為「DRAM 搬運速度用滿」。
分散的存取增加每個請求涉及的 sector 數，
與記憶體指令佇列壓力及低 FP32 效能的觀察一致。

這也不違反 Part 0：Part 0 的 compute-bound 判斷建立在 A、B 各從 DRAM
讀取一次、C 寫回一次的理想假設，並未保證 naive 實作能有效利用硬體。
效能比較採用 benchmark 時間；profiler 指標則解釋存取行為與等待原因，
兩者的時間不混用。

\clearpage
\section*{Part 2: Coalesced Global Memory Access}

\subsection*{Q1. Benchmark 結果}

Coalesced 與 naive 皆由每個 thread 計算一個 C 元素，
依序沿著 $k$ 累加 A 與 B 的乘積，最後計算
$C=\alpha AB+\beta C$。Coalesced 的 thread mapping 為
\texttt{threadIdx.x} 對應 col、\texttt{threadIdx.y} 對應 row，
Block 為 $(32,32)$，grid 為
$(\lceil N/32\rceil,\lceil M/32\rceil)$，邊界檢查維持
\texttt{row < M \&\& col < N}，因此也能處理非整除與非方形矩陣。
此版本不使用 shared memory，效能差異來自同一 warp 內的位址排列。

在 $\alpha=1$、$\beta=0$ 下，三種尺寸的結果皆為 \texttt{PASS}，
相對 naive 約加速 $5.59$--$10.15\times$。
加速倍率為相同尺寸的 naive 時間除以 coalesced 時間。

\begin{table}[H]
\centering
\small
\begin{tabular}{@{}rrrrrcr@{}}
\toprule
$M=N=K$ & ms & GFLOP/s & \%cuBLAS & \%peakF & Check & 相對 naive \\
\midrule
1024 & 6.3365   & 338.9 & 13.1 & 4.2 & PASS & $5.59\times$ \\
2048 & 27.4532  & 625.8 & 15.3 & 7.7 & PASS & $10.15\times$ \\
4096 & 223.6519 & 614.5 & 14.6 & 7.5 & PASS & $9.98\times$ \\
\bottomrule
\end{tabular}
\caption{Coalesced 的一般 benchmark 結果；三種尺寸皆通過正確性檢查。}
\label{tab:part2-benchmark}
\end{table}

\reportscreenshot{graph/part2_1024.png}
  {$1024^3$ 的 coalesced benchmark 原始輸出。}
  {fig:part2-benchmark-1024}
\reportscreenshot{graph/part2_2048.png}
  {$2048^3$ 的 coalesced benchmark 原始輸出。}
  {fig:part2-benchmark-2048}
\reportscreenshot{graph/part2_4096.png}
  {$4096^3$ 的 coalesced benchmark 原始輸出。}
  {fig:part2-benchmark-4096}

\textbf{正確性與 sanitizer 檢查。}
Coalesced 在 15 種尺寸、$\alpha=1.5$、$\beta\in\{0,0.5\}$
的 30 組正確性測試皆為 \texttt{PASS}。
5 組 memcheck 皆為 0 errors；3 組 racecheck 皆為 0 errors、0 warnings。

\clearpage
\subsection*{Q2. 相同運算量與讀取次數，為何速度改變？}

\textbf{量測結果。}
在 $2048^3$ 下，兩者的 global-load requests 數量相同，
但 coalesced 的 sectors/request 明顯較低。
這裡的 request 是 warp 記憶體指令送往 L1/TEX 的要求，
不是一次 DRAM 傳輸。

\begin{table}[H]
\centering
\small
\begin{tabular}{@{}lrr@{}}
\toprule
L1/TEX 指標 & Naive & Coalesced \\
\midrule
Global-load requests & $537{,}001{,}984$ & $537{,}001{,}984$ \\
Global-load sectors & $8{,}862{,}535{,}049$ & $1{,}342{,}701{,}568$ \\
Load sectors/request & $16.50$ & $2.50$ \\
Global-store requests & $131{,}072$ & $131{,}072$ \\
Global-store sectors & $4{,}194{,}304$ & $524{,}288$ \\
Store sectors/request & $32.00$ & $4.00$ \\
\bottomrule
\end{tabular}
\caption{相同尺寸下，coalescing 前後的 L1/TEX 存取統計。}
\label{tab:part2-sectors}
\end{table}

Coalesced 的 global-load sectors/request 為
\[
\frac{1{,}342{,}701{,}568}{537{,}001{,}984}
\approx2.50,
\]
相較 naive 的 $16.50$ 明顯降低。
兩者 global-load requests 完全相同，但 coalesced 的 sector 存取總數
減少約 $84.85\%$。這是 L1/TEX 的 sector 存取量改善，
\emph{不能解讀成 DRAM 流量或執行時間也減少相同比例}。

\textbf{由 thread mapping 解釋。}
在主要迴圈固定的 $k$ 下，naive 的同一個 warp 負責不同 row、相同 col，
讀 A 涉及 32 個 sector，讀 B 則只涉及 1 個 sector，
因此兩次讀取平均為 $(32+1)/2=16.5$。

Coalesced 改為同一個 warp 負責相同 row、連續的 32 個 col。
此時 32 個 thread 的 \texttt{A[row*K+k]} 是同一個元素，只涉及 1 個 sector；
\texttt{B[k*N+col]} 則為連續的 32 個 FP32 元素。
一個 sector 為 32 bytes，可容納 8 個 FP32 元素，因此對本次對齊的
$2048^3$ 存取，只需 4 個 sector 服務這次 B 讀取。
主要迴圈的平均為
\[
\frac{1+4}{2}=2.5\ \text{sectors/request},
\]
與實測相符；完整統計還包含迴圈外對 C 的讀取，因此精確比值略高於 $2.5$。
寫回 C 時，warp 內的位址也由跨列分散變為連續，
所以 store sectors/request 從 $32$ 降為 $4$。

每個 thread 仍執行相同次數的乘加，並以相同的 $k$ 順序讀取自己的 A、B 資料；
改變的是\emph{同一時間、同一個 warp 內不同 thread 的位址關係}。
更集中的存取降低記憶體請求的內部處理工作，
不需要減少 FLOP 或讀取指令數，就能改善效能。

\reportscreenshot{graph/part2_global.png}
  {Coalesced 的 L1/TEX global load/store requests 與 sectors 原始計數。}
  {fig:part2-global-access}

\clearpage
\subsection*{Q3. 哪個記憶體層級最忙？重複讀取由哪裡服務？}

\textbf{記憶體階層的吞吐。}
Coalesced 的 L1/TEX Cache Throughput 為
$83.67\%$，高於 L2 的 $6.95\%$ 與 DRAM 的 $24.25\%$。
以這些各自正規化的吞吐指標判斷，L1/TEX 最接近其處理上限，
而不是 DRAM 頻寬已經飽和。
這些百分比的計數方式與 active/elapsed 基準不完全相同，
不能相加或當成資料來自各階層的分配比例。

\begin{table}[H]
\centering
\begin{tabular}{@{}lrr@{}}
\toprule
GPU Speed Of Light 指標 & Naive & Coalesced \\
\midrule
Duration & $278.23\ \mathrm{ms}$ & $20.49\ \mathrm{ms}$ \\
Compute (SM) Throughput & $6.03\%$ & $83.66\%$ \\
Memory Throughput & $49.77\%$ & $83.66\%$ \\
L1/TEX Cache Throughput & $99.53\%$ & $83.67\%$ \\
L2 Cache Throughput & $0.70\%$ & $6.95\%$ \\
DRAM Throughput & $2.49\%$ & $24.25\%$ \\
\bottomrule
\end{tabular}
\caption{Coalescing 前後的 GPU Speed Of Light 指標。Duration 為 profiling 量測。}
\label{tab:part2-sol}
\end{table}

\reportscreenshot{graph/part2_sol.png}
  {Coalesced 的 GPU Speed Of Light Throughput 原始輸出。}
  {fig:part2-sol}

\textbf{重複讀取主要由 cache 服務。}
雖然從 thread 的計算需求來看，每個 A 元素會被不同 col 的 thread
讀取 $N$ 次，每個 B 元素會被不同 row 的 thread 讀取 $M$ 次，
這並不代表每次讀取都要到 DRAM。
Coalesced 的 L1/TEX sector 命中率為 $94.96\%$，
顯示大量 sector 存取在 L1 就能命中；部分未命中的存取再由 L2 服務。
因此，大量重複讀取主要由 L1 cache 吸收，而非一律重新從 DRAM 搬入。

\begin{table}[H]
\centering
\begin{tabular}{@{}lrr@{}}
\toprule
Memory Workload Analysis 指標 & Naive & Coalesced \\
\midrule
Memory Throughput（DRAM 讀寫合計）& $8.00\ \mathrm{GB/s}$ & $83.15\ \mathrm{GB/s}$ \\
L1/TEX Hit Rate & $99.04\%$ & $94.96\%$ \\
L2 Hit Rate & $82.07\%$ & $49.62\%$ \\
Mem Pipes Busy & $6.03\%$ & $83.66\%$ \\
\bottomrule
\end{tabular}
\caption{Coalescing 前後的實際 DRAM 吞吐、cache 命中率與記憶體指令吞吐指標。}
\label{tab:part2-memory}
\end{table}

L1 與 L2 命中率各自以到達該層級的 sector 查詢為分母，不能直接相加。
Coalesced 的命中率雖較低，執行仍更快，因為命中率只描述資料是否在 cache，
並未描述每個 request 涉及多少 sector。
Naive 即使有很高的命中率，仍需處理大量分散的 sector 存取。
另外，此處 $83.15\ \mathrm{GB/s}$ 是 profiler 量到的實際 DRAM 讀寫吞吐，
不同於 harness 用最低必要資料量除以時間得到的 \texttt{GB/s}。

\reportscreenshot{graph/part2_memory.png}
  {Coalesced 的 Memory Workload Analysis 原始輸出。}
  {fig:part2-memory}

\textbf{Compute (SM) 高，不等於 FP32 乘加已飽和。}
Compute Workload Analysis 的
\emph{Pipe Utilization (\% of peak instructions executed)} 表中，
LSU 為 $83.44\%$，FMA 為 $12.46\%$。
LSU 負責發出 load/store 等操作，FMA 管線則主要處理浮點算術，
也可能執行部分整數運算。
因此，Compute (SM) 的高吞吐不能直接當成 FP32 算力的使用比例；
LSU 的高利用率與記憶體指令處理管線仍承受較大壓力的判斷一致。

實際 FP32 效能仍以一般 benchmark 為準。以 $2048^3$ 為例：
\[
\frac{625.8}{8140.8}\times100\%\approx7.69\%,
\]
與 harness 的 $7.7\%$ 相符，而不是 $83.66\%$。

\reportscreenshot{graph/part2_compute.png}
  {Compute Workload Analysis 的管線利用率，包含 LSU 與 FMA 的實測值。}
  {fig:part2-compute}

\textbf{LG Throttle 明顯改善，但仍是主要等待原因。}
LG Throttle 由 naive 的 $365.25$ 降至 coalesced 的 $20.38$。
Warp State Statistics 的數值是每個 issued instruction 對應的平均 warp cycles，
不是整張 GPU 完全閒置的時間。

\begin{table}[H]
\centering
\begin{tabular}{@{}lrr@{}}
\toprule
平均週期／issued instruction & Naive & Coalesced \\
\midrule
Warp Cycles Per Issued Instruction & $373.61$ & $27.25$ \\
Stall LG Throttle & $365.25$ & $20.38$ \\
Stall Long Scoreboard & $3.41$ & $1.27$ \\
\bottomrule
\end{tabular}
\caption{Coalescing 前後的 warp 指令間隔與主要記憶體等待指標。}
\label{tab:part2-stalls}
\end{table}

Coalesced 的 LG Throttle 約占平均指令間隔的
\[
\frac{20.38}{27.25}\times100\%\approx74.79\%.
\]
雖然仍是最大的等待原因，但其絕對值已從 $365.25$ 降為 $20.38$。
這與 sector/request 降低、記憶體要求更容易處理的結果一致。
LG Throttle 仍存在也合理：主要迴圈仍持續發出大量 global loads，
Part 2 只改善存取集中程度，尚未透過 block 內明確共享來減少這些重複讀取。
Profiler 的 Estimated Speedup 僅為潛在改善估計，
不當成已測得或保證可達成的加速比例。

\reportscreenshot[trim=75bp 165bp 188bp 182.5bp,clip]{graph/part2_warp.png}
  {Coalesced 的 Warp State Statistics 節錄：平均 warp cycles 與主要等待指標。}
  {fig:part2-warp}

\textbf{結論。}
Coalescing 在不改變主要 FLOP 與 global-load requests 數量的情況下，
將平均 load sectors/request 從 $16.50$ 降至 $2.50$，
並伴隨 LG Throttle 等待大幅減少。
三種尺寸的一般 benchmark 相對 naive 約加速 $5.59$--$10.15\times$。
重複讀取大量由 L1 cache 服務；剩餘壓力主要指向記憶體存取的處理，
不能簡化為 FP32 算力或 DRAM 頻寬已經用滿。

\clearpage
\begingroup
\setlength{\emergencystretch}{2em}
\section*{Part 3: Shared Memory and Register Blocking}

\subsection*{Q1. Benchmark 與六組參數實驗}

在 $M=N=K=4096$、$\alpha=1$、$\beta=0$ 下，
smem16、smem32 與六組 register-blocking 的結果皆為 \texttt{PASS}。
六組參數以 R1--R6 表示，涵蓋不同的 block tile、K tile 與 thread tile 大小。
Register blocking 的效能為 $2004.5$--$3157.4\ \mathrm{GFLOP/s}$，
均高於每個 thread 只計算一個輸出的 smem16 與 smem32。

\begin{table}[H]
\centering
\small
\begin{tabular}{@{}lrrrrc@{}}
\toprule
Kernel & ms & GFLOP/s & \%cuBLAS & \%peakF & Check \\
\midrule
smem16 & 207.5528 & 662.2  & 15.7 & 8.1  & PASS \\
smem32 & 158.0202 & 869.8  & 20.7 & 10.7 & PASS \\
R1     & 45.6292  & 3012.1 & 72.6 & 37.0 & PASS \\
R2     & 68.5650  & 2004.5 & 48.3 & 24.6 & PASS \\
R3     & 52.0103  & 2642.5 & 63.6 & 32.5 & PASS \\
R4     & 47.8910  & 2869.8 & 69.1 & 35.3 & PASS \\
R5     & 43.5296  & 3157.4 & 75.1 & 38.8 & PASS \\
R6     & 50.1503  & 2740.5 & 65.3 & 33.7 & PASS \\
\bottomrule
\end{tabular}
\caption{$4096^3$ 的 benchmark 結果。}
\label{tab:part3-benchmark}
\end{table}

\begin{table}[H]
\centering
\small
\begin{tabular}{@{}lccrrr@{}}
\toprule
設定 & $BM\times BN\times BK$ & $TM\times TN$ & Threads/block
     & Regs/thread & Smem/block (B) \\
\midrule
R1 & $128\times128\times8$  & $8\times8$ & 256 & 119 & 8192 \\
R2 & $64\times64\times8$    & $4\times4$ & 256 & 62  & 4096 \\
R3 & $128\times64\times8$   & $8\times4$ & 256 & 76  & 6144 \\
R4 & $128\times128\times16$ & $8\times8$ & 256 & 116 & 16384 \\
R5 & $128\times128\times8$  & $8\times4$ & 512 & 64  & 8192 \\
R6 & $64\times128\times8$   & $8\times8$ & 128 & 117 & 6144 \\
\bottomrule
\end{tabular}
\caption{六組參數及編譯器資源使用量；六組皆無 register spill stores/loads。}
\label{tab:part3-resources}
\end{table}

Registers/thread 為 ptxas 回報的每個 thread 的 register 配置量。
Shared memory/block 則由 A、B 兩個 FP32 tile 的大小決定：
\[
S=4(BM\cdot BK+BK\cdot BN)=4BK(BM+BN)\ \text{bytes}.
\]
例如 R1 為 $4\times8\times(128+128)=8192$ bytes。
\clearpage
\textbf{Occupancy 與主要 warp stall。}
理論 occupancy 為資源配置允許的 resident warps 占硬體上限的比例，
實測 occupancy 則反映執行期間的平均 active warps。
主要 stall 的數值為每個 issued instruction 對應的平均 warp cycles，
不是百分比或 kernel 的總等待時間。

\begin{table}[H]
\centering
\small
\begin{tabular}{@{}lrrlr@{}}
\toprule
設定 & 理論 occupancy & 實測 occupancy & 最大 stall & 平均週期 \\
\midrule
R1 & $50\%$  & $49.25\%$ & Short Scoreboard & 1.64 \\
R2 & $100\%$ & $98.99\%$ & Long Scoreboard  & 7.45 \\
R3 & $75\%$  & $73.87\%$ & Long Scoreboard  & 3.51 \\
R4 & $50\%$  & $49.29\%$ & MIO Throttle     & 2.52 \\
R5 & $100\%$ & $98.36\%$ & Long Scoreboard  & 4.33 \\
R6 & $50\%$  & $49.00\%$ & Long Scoreboard  & 2.09 \\
\bottomrule
\end{tabular}
\caption{六組 register-blocking 的 occupancy 與最大等待原因。}
\label{tab:part3-occupancy}
\end{table}

\begin{table}[H]
\centering
\small
\begin{tabular}{@{}lrrrr@{}}
\toprule
設定 & Shared-load 指令 & Shared-store 指令 & Load conflicts & Store conflicts \\
\midrule
R1 & $134{,}217{,}728$ & $33{,}554{,}432$ & $268{,}435{,}456$ & $117{,}440{,}512$ \\
R2 & $268{,}435{,}456$ & $67{,}108{,}864$ & $0$             & $234{,}881{,}024$ \\
R3 & $201{,}326{,}592$ & $50{,}331{,}648$ & $0$             & $234{,}881{,}024$ \\
R4 & $134{,}217{,}728$ & $33{,}554{,}432$ & $268{,}435{,}456$ & $251{,}658{,}240$ \\
R5 & $201{,}326{,}592$ & $33{,}554{,}432$ & $0$             & $117{,}440{,}512$ \\
R6 & $134{,}217{,}728$ & $50{,}331{,}648$ & $268{,}435{,}456$ & $117{,}440{,}512$ \\
\bottomrule
\end{tabular}
\caption{Shared Memory 表的 SASS 指令與 shared load/store bank-conflict 原始計數。}
\label{tab:part3-shared}
\end{table}

這些是整個 kernel 的計數，不是每個 thread 的次數。
Bank conflict 表示同一 shared-memory bank 的不同位址存取需要額外處理，
不等於資料競爭或計算錯誤。讀取同一個位址可 broadcast，不能僅由
「多個 thread 存取同一個 bank」判定一定有 conflict。
寬指令本來也可能需要多個 wavefront，因此不直接把
wavefronts/instruction 當成全部可消除的衝突倍率。

\textbf{正確性。}
完整正確性與 sanitizer 檢查結果為 \texttt{ALL PASSED}。
\texttt{naive}、\texttt{coalesced}、\texttt{smem16}、\texttt{smem32}
各在 15 種尺寸、兩種 $\beta$ 的 30 組測試皆為 \texttt{PASS}。
R2 為 10 組 \texttt{PASS}、20 組 \texttt{SKIPPED}，其餘五組各為
8 組 \texttt{PASS}、22 組 \texttt{SKIPPED}，皆符合 launcher 的尺寸限制。
5 組 memcheck 與 3 組 racecheck 也皆為 clean。
其中 \texttt{SKIPPED} 表示尺寸不符合 register-blocking 的整除條件，
不代表該尺寸的計算結果已通過檢查。

\reportscreenshot{graph/part3a_smem16_4096.png}
  {$4096^3$ 的 smem16 benchmark。}
  {fig:part3a-smem16}
\reportscreenshot{graph/part3a_smem32_4096.png}
  {$4096^3$ 的 smem32 benchmark。}
  {fig:part3a-smem32}

\clearpage
\subsection*{Q2. 每次 FMA 需要多少 shared-memory loads？}

只計算內層乘加迴圈在原始碼層級讀取的 FP32 元素數，
不把 global-to-shared 載入、shared stores 或 C 的寫回混入此比值。
Part 3a 每個 thread 計算一個輸出，每個 $k$ 讀取
\texttt{As[ty][k]} 與 \texttt{Bs[k][tx]}，完成一次 FMA，因此
\[
L_{\mathrm{3a}}=\frac{1+1}{1}=2\ \text{shared loads/FMA}.
\]

Part 3b 每個 $k$ 先將 $TM$ 個 A 與 $TN$ 個 B 元素讀入 registers，
再將這些值交叉組合，完成 $TM\times TN$ 次 FMA：
\[
L_{\mathrm{3b}}=\frac{TM+TN}{TM\cdot TN}.
\]
當 $TM=TN=8$，16 個 shared-memory 元素供 64 次 FMA 重複使用，故
\[
L_{\mathrm{3b},8\times8}=\frac{8+8}{8\times8}=0.25,
\qquad
\frac{L_{\mathrm{3a}}}{L_{\mathrm{3b},8\times8}}=8.
\]
也就是每次 FMA 的 shared-memory 元素讀取需求降為原本的八分之一。
這是 \emph{register reuse}：A 值可供多個輸出 column 使用，
B 值可供多個輸出 row 使用，不是省略矩陣乘法需要的乘加。

以同為 $8\times8$ thread tile 的 R1 為例，一般 benchmark 的時間為
$45.6292\ \mathrm{ms}$，相對 smem16 與 smem32 分別約快
\[
\frac{207.5528}{45.6292}\approx4.55\times,
\qquad
\frac{158.0202}{45.6292}\approx3.46\times.
\]
較少的 shared-memory 元素讀取需求有利於降低記憶體指令壓力，
但不保證得到 $8\times$ 的執行加速；global 載入、同步、bank conflicts、
register 用量與 occupancy 仍會影響時間。
此外，編譯器可能合併成寬指令，所以此處的「元素 loads/FMA」
不等於實際的 SASS 指令數/FMA。

\subsection*{Q3. 最快與較慢設定的差異}

\textbf{R5 與 R2：相近的高 occupancy，不同的指令工作量與等待。}
R5 的量測時間最短，R2 最長，兩者的實測加速為
\[
\frac{68.5650}{43.5296}\approx1.58\times.
\]
兩者的理論 occupancy 都是 $100\%$，實測分別為 $98.36\%$ 與 $98.99\%$，
因此不能將 R5 較快歸因於 occupancy 更高。
R2 每個 thread 計算 $4\times4$，每次 FMA 需
$(4+4)/16=0.5$ 個 shared-memory 元素 loads；
R5 每個 thread 計算 $8\times4$，則為 $(8+4)/32=0.375$，減少 $25\%$。
R5 的 $128\times128$ block tile 也比 R2 的 $64\times64$ 大，
能在 block 內重用更多資料。

實測 shared-load SASS 指令由 R2 的 $268{,}435{,}456$ 降至
R5 的 $201{,}326{,}592$，shared-store 指令則由 $67{,}108{,}864$
降至 $33{,}554{,}432$。兩者 load bank conflicts 都是 0；
R5 的 store conflicts 雖減半，但 store 指令數也減半，
不能將其解釋成每次 store 的衝突程度已改善。
因此，即使 load 存取沒有 bank conflicts，整體指令工作量仍能造成效能差距。

Warp 指標亦支持 R2 有更大的資料依賴與指令佇列壓力：
Long Scoreboard 從 $7.45$ 降至 R5 的 $4.33$，
MIO Throttle 從 $6.05$ 降至 $1.61$，Barrier 從 $2.44$ 降至 $1.53$。
Long Scoreboard 是等待 L1/TEX 操作的結果，不代表一定等 DRAM；
MIO Throttle 是相關指令佇列壓力，Barrier 則是等同 block 的其他 warps。
這些平均 warp 週期不能直接相加換算成時間加速。
R5 的 Short Scoreboard 反而較高（$3.55$ 對 $1.47$），
顯示最快設定也不必在每一項 stall 上都最小。

\textbf{其餘參數的取捨。}
R1 的 $8\times8$ thread tile 有較佳的 register reuse，
但每個 thread 使用 119 個 registers，register 限制每個 SM 只能放 2 個
256-thread blocks，即 $2\times8=16$ 個 warps，理論 occupancy 為 $50\%$。
R5 改為 $8\times4$ 後，每個 thread 使用 64 個 registers；
每個 block 有 512 threads，仍可放 2 個 blocks，卻達到 32 個 resident warps。
R5 的 load conflicts 也由 R1 的 $268{,}435{,}456$ 降為 0。
這與 R5 較快的結果一致，但兩組同時改變多項執行特性，
不能將全部差異歸給單一因素。

R4 將 R1 的 $BK$ 由 8 增至 16，外層 tile 迭代數減半，
Barrier 由 $0.80$ 降至 $0.61$；然而 shared memory/block 加倍，
且 store bank conflicts 從 $117{,}440{,}512$ 增至 $251{,}658{,}240$，
MIO Throttle 亦從 $1.44$ 增至 $2.52$，執行時間反而較長。
R6 與 R1 同為 $8\times8$ thread tile，僅將 $BM$ 減半；
兩者理論 occupancy 都是 $50\%$，load/store conflict 總數也相同，
但 R6 的 shared-store 指令由 $33{,}554{,}432$ 增至 $50{,}331{,}648$，
時間增至 $50.1503\ \mathrm{ms}$。
這再次顯示 occupancy 或 conflict 總數本身不足以預測效能。

綜合而言，較大的 register tile 可提高資料重用率，
但也可能增加 register 用量、降低 occupancy 或改變 bank conflicts。
最佳參數取決於這些因素的平衡，不能只以單一指標判斷。

\clearpage
\subsection*{六組 Benchmark 與 Profiler 截圖}

R1--R6 的 benchmark、Occupancy、Warp State 與 Shared Memory 量測結果如下。

% Display-only crops remove terminal margins; the source PNGs are unchanged.
\newcommand{\regprofileimage}[2][]{%
  \IfFileExists{#2}{%
    \includegraphics[#1,width=\linewidth]{#2}%
  }{%
    \fbox{\parbox{0.92\linewidth}{\centering\small
      圖片尚未載入\\\path{#2}}}%
  }%
}
\newcommand{\regprofileevidence}[4]{%
  \begin{figure}[H]
    \centering
    {\small\textbf{#1：}\texttt{\detokenize{#2}}}\par\smallskip
    \regprofileimage{graph/part3b_#2_4096.png}\par\medskip
    \begin{minipage}[t]{0.47\linewidth}
      \vspace{0pt}\centering
      {\small (a) Occupancy}\par\smallskip
      \regprofileimage[#3]{graph/part3b_#2_occupancy_4096.png}
    \end{minipage}\hfill
    \begin{minipage}[t]{0.47\linewidth}
      \vspace{0pt}\centering
      {\small (b) Warp State}\par\smallskip
      \regprofileimage[#4]{graph/part3b_#2_warp_4096.png}
    \end{minipage}\par\medskip
    {\small (c) Shared Memory}\par\smallskip
    \regprofileimage{graph/part3b_#2_shared_4096.png}
    \caption{#1 在 $4096^3$ 的 benchmark 與分析指標原始截圖。}
    \label{fig:part3b-#1}
  \end{figure}
}

\regprofileevidence{R1}{reg2d_128x128x8_8x8}
  {trim=22.5bp 60.75bp 752.25bp 85.5bp,clip}
  {trim=22.5bp 6bp 843.75bp 178.5bp,clip}
\clearpage
\regprofileevidence{R2}{reg2d_64x64x8_4x4}
  {trim=15bp 23.25bp 599.25bp 46.5bp,clip}
  {trim=0bp 0bp 0bp 0bp,clip}
\clearpage
\regprofileevidence{R3}{reg2d_128x64x8_8x4}
  {trim=15bp 39.75bp 599.25bp 46.5bp,clip}
  {trim=0bp 0bp 0bp 0bp,clip}
\clearpage
\regprofileevidence{R4}{reg2d_128x128x16_8x8}
  {trim=15bp 50.25bp 609bp 46.5bp,clip}
  {trim=0bp 0bp 0bp 0bp,clip}
\clearpage
\regprofileevidence{R5}{reg2d_128x128x8_8x4}
  {trim=15bp 23.25bp 609bp 46.5bp,clip}
  {trim=0bp 0bp 0bp 0bp,clip}
\clearpage
\regprofileevidence{R6}{reg2d_64x128x8_8x8}
  {trim=15bp 42.75bp 603bp 46.5bp,clip}
  {trim=0bp 0bp 0bp 0bp,clip}

\clearpage
\endgroup

\begingroup
\section*{Part 4: Further Optimization}

兩項優化皆以 Part 3 的 R5（\texttt{reg2d\_128x128x8\_8x4}）為起點，
維持 $BM=BN=128$、$BK=8$、$TM=8$、$TN=4$，每個 block 有 512 個 threads。
\texttt{custom\_pad4} 只調整 A tile 的 shared-memory 排列；
\texttt{custom\_pad4\_vec4} 保留此排列，再將 A、B 的 global loads 向量化。
兩者均維持 FP32 乘加與 $C=\alpha AB+\beta C$ 的計算，不使用 tensor cores。
以下效能比較皆為 $M=N=K=4096$、$\alpha=1$、$\beta=0$ 的最終評分量測。

% Display-only crops retain the metric rows used in the analysis.
\newcommand{\partfourimage}[2][]{%
  \IfFileExists{#2}{%
    \includegraphics[#1,width=\linewidth]{#2}%
  }{%
    \fbox{\parbox{0.92\linewidth}{\centering\small
      圖片尚未載入\\\path{#2}}}%
  }%
}

\subsection*{優化一：Shared-memory Padding}

\textbf{設計變更。}
A tile 的宣告由 \texttt{As[BK][BM]} 改為 \texttt{As[BK][BM+4]}，
其餘載入迴圈、thread mapping、內層乘加與 C 寫回皆不變。
每列仍只使用前 $BM$ 個元素；額外四個位置不參與計算，
而是將相鄰列的儲存間距由 128 個 float 改為 132 個 float。

\textbf{為何能減少 bank conflicts？}
Shared memory 的連續 FP32 元素依序對應 32 個 banks。
以 \texttt{As[0][0]} 位於 bank 0 為例，原本與 padding 後的 bank 編號分別為
\[
\begin{aligned}
b_{\mathrm{original}}(k,m)&=(128k+m)\bmod32=m\bmod32,\\
b_{\mathrm{pad4}}(k,m)&=(132k+m)\bmod32=(4k+m)\bmod32.
\end{aligned}
\]
在 cooperative loading 的一次迭代中，同一 warp 的
\texttt{idx} 為連續 32 個值，且 $m=\lfloor\texttt{idx}/8\rfloor$、
$k=\texttt{idx}\bmod8$。
因此原本每八個 threads 寫入同一 bank 的八個不同位置，
整個 warp 僅分布在四個 banks；padding 後則分散至 32 個不同 banks。
這改善的是轉置寫入 \texttt{As[k][m]} 的衝突，不是減少矩陣元素或 FLOP。

\begin{table}[H]
\centering
\small
\begin{tabular}{@{}lrr@{}}
\toprule
Profiler 指標 & R5 原版 & \texttt{custom\_pad4} \\
\midrule
Shared-load bank conflicts & $0$ & $0$ \\
Shared-store bank conflicts & $117{,}440{,}512$ & $0$ \\
\bottomrule
\end{tabular}
\caption{$4096^3$ 的 shared-memory bank-conflict 原始計數。}
\end{table}

\textbf{效能與代價。}
時間由 $45.959\ \mathrm{ms}$ 降至 $42.117\ \mathrm{ms}$，
相對原版加速 $1.091\times$。
Shared-store bank conflicts 降為 0，與減少存取序列化的設計預期一致。
Shared memory/block 由 8192 bytes 增至 8320 bytes，僅多
$BK\times4\times4=128$ bytes；兩版均為 64 registers/thread，且無 register spills。
因此本次改善不需改變 thread tile 或增加乘加運算，代價主要是少量額外儲存空間。

\begin{figure}[H]
\centering
{\small\textbf{R5 原版：shared-memory bank conflicts}}\par\smallskip
\partfourimage[trim=0bp 120bp 0bp 0bp,clip]
  {graph/part3b_reg2d_128x128x8_8x4_shared_4096.png}\par\medskip
{\small\textbf{\texttt{custom\_pad4}：shared-memory bank conflicts}}\par\smallskip
\partfourimage[trim=0bp 120bp 0bp 0bp,clip]
  {graph/part4_custom_pad4_shared_4096.png}
\caption{Padding 前後的 shared-memory 指標節錄，包含 load 與 store bank conflicts。}
\end{figure}

\clearpage
\subsection*{優化二：Vectorized Global Loads}

\textbf{設計變更。}
\texttt{custom\_pad4\_vec4} 保留 padding 與既有乘加迴圈，
僅將 A、B tile 的 scalar global loads 改為 \texttt{float4} loads。
原本每次處理一個 FP32 元素，現在每次讀取四個連續 FP32 元素，
再以 \texttt{v.x}、\texttt{v.y}、\texttt{v.z}、\texttt{v.w}
分別放入對應的 shared-memory 位置。
A 仍以 \texttt{As[k][m]} 轉置儲存，B 仍以 \texttt{Bs[k][n]} 儲存。
載入迴圈的工作單位改為四個元素，因此 A、B 的迴圈上限分別為
$BM(BK/4)$ 與 $BK(BN/4)$，並相應調整索引。

\textbf{為何能改善效能？}
\texttt{float4} 將四個 float 組成 16-byte 的讀取單位，
可用較少的 global-load 指令取得相同資料。
本設定的 $BK=8$、$BN=128$ 及 launcher 的整除限制，配合對齊的矩陣起始位址，
使向量化讀取符合 16-byte 對齊要求。
乘加次數、thread 數量、C 的寫回與 FP32 精度皆維持不變。

\begin{table}[H]
\centering
\small
\begin{tabular}{@{}lrr@{}}
\toprule
Profiler 指標 & \texttt{custom\_pad4} & \texttt{custom\_pad4\_vec4} \\
\midrule
Global-load SASS 指令 & $34{,}078{,}720$ & $8{,}912{,}896$ \\
Global-store SASS 指令 & $524{,}288$ & $524{,}288$ \\
\bottomrule
\end{tabular}
\caption{向量化前後的動態指令計數，分別對應 global load/store 的 SASS 指標。}
\end{table}

Global-load 指令數減少約
$\bigl(1-8{,}912{,}896/34{,}078{,}720\bigr)\times100\%=73.85\%$，
而 global-store 指令數不變，與僅修改載入方式的設計相符。
這些是整個 kernel 的指令計數，並非讀取的 float 數或 DRAM bytes；
不能據此認為 DRAM 流量也減少 $73.85\%$，或時間應縮短四倍。
時間由 $42.117\ \mathrm{ms}$ 降至 $36.235\ \mathrm{ms}$，
相對 padding 版本加速 $1.162\times$。
兩版均使用 64 registers/thread、8320 bytes shared memory/block，且無 register spills。

\begin{figure}[H]
\centering
{\small\textbf{\texttt{custom\_pad4}：global-load/store 指令}}\par\smallskip
\partfourimage[trim=0bp 1.5bp 363bp 36bp,clip]
  {graph/part4_custom_pad4_global_4096.png}\par\medskip
{\small\textbf{\texttt{custom\_pad4\_vec4}：global-load/store 指令}}\par\smallskip
\partfourimage[trim=0bp 1.5bp 393bp 36bp,clip]
  {graph/part4_custom_pad4_vec4_global_4096.png}
\caption{向量化前後的 global-load/store SASS 指令原始計數節錄。}
\end{figure}

\subsection*{效能比較與正確性}

最終評分量測對每個版本取三輪、每輪 20 次執行的中位數，再取三輪中的最佳值。
加速倍率為前一版本的時間除以目前版本時間。
兩步皆超過 $1.05\times$ 的門檻，各得 5 分，Part 4 合計為 $10/10$。
下列時間與前述兩項優化的比較來自同一次評分結果，
不混用其他 benchmark 或 profiler 的時間。

\begin{table}[H]
\centering
\small
\begin{tabular}{@{}lrrc@{}}
\toprule
Kernel & ms & 相對前一步 & 得分 \\
\midrule
\texttt{reg2d\_128x128x8\_8x4} & 45.959 & 基準 & -- \\
\texttt{custom\_pad4} & 42.117 & $1.091\times$ & 5 \\
\texttt{custom\_pad4\_vec4} & 36.235 & $1.162\times$ & 5 \\
\bottomrule
\end{tabular}
\caption{$4096^3$ 的最終優化鏈與評分結果。}
\end{table}

兩項自訂 kernel 均在支援尺寸通過正確性、memcheck 與 racecheck 檢查。
不符合 launcher 整除條件的尺寸為 \texttt{SKIPPED}，不當成已通過的計算結果。
因此兩項改善均建立在正確性檢查通過的基礎上：
padding 改善 shared-memory 寫入的 bank conflicts，
向量化則降低 global-load 指令工作量。

\reportscreenshot[trim=0bp 0bp 0bp 277.5bp,clip]{graph/part4_grade_vec4.png}
  {最終評分結果節錄：兩項優化皆成功，Part 4 為 $10/10$。}
  {fig:part4-grade}
\endgroup
\end{document}
